\section{Приложение: справочник платформы}
\subsection{Ключевые события \texttt{Ev.*}}
\begin{itemize}
  \item \texttt{Ev.MCE\_PREFLIGHT} — безопасная предстартовая подготовка.
  \item \texttt{Ev.MCE\_TAKEOFF} — команда взлёта; завершение сигналом \texttt{Ev.TAKEOFF\_COMPLETE}.
  \item \texttt{Ev.MCE\_LANDING} — команда посадки; завершение \texttt{Ev.COPTER\_LANDED}.
  \item \texttt{Ev.ENGINES\_DISARM} — выключение двигателей.
  \item \texttt{Ev.POINT\_REACHED} — целевая точка достигнута.
  \item \texttt{Ev.LOW\_VOLTAGE2} — низкое напряжение (аварийная реакция).
  \item \texttt{Ev.SHOCK} — удар/столкновение (экстренная остановка).
\end{itemize}

\subsection{Основные функции управления}
\begin{itemize}
  \item Команды смены режима (через \texttt{ap.push}):
  \begin{itemize}
      \item \texttt{Ev.MCE\_PREFLIGHT} — предстартовая подготовка.
      \item \texttt{Ev.MCE\_TAKEOFF} — взлёт.
      \item \texttt{Ev.MCE\_LANDING} — посадка.
      \item \texttt{Ev.ENGINES\_DISARM} — выключение двигателей.
  \end{itemize}
  \item \texttt{ap.goToLocalPoint(x,y,z)} — полёт к точке в локальной системе координат (метры).
  \item \texttt{ap.updateYaw(rad)} — установка курса в радианах; перевод \texttt{deg} в радианы: \\ \texttt{rad = deg * math.pi / 180}.
  \item \texttt{Sensors.rc()} — чтение RC‑каналов; проверка тумблера: \texttt{rc\_chans[8]}.
\end{itemize}

\subsection{Таймеры и неблокирующие вызовы}
\begin{itemize}
  \item \texttt{Timer.new(period, fn)} — периодический вызов \texttt{fn} каждые \texttt{period} секунд.
  \item \texttt{Timer.callLater(delay, fn)} — отложенный однократный вызов без блокировки.
  \item \texttt{start()/stop()} — управление таймером; избегайте \texttt{sleep()} в миссиях.
\end{itemize}

\subsection{Идиомы Lua}
\begin{itemize}
\item Распаковка: \texttt{local unpack = table.unpack}; передача RGB: \texttt{leds:set(i, unpack(col))}.
\item «Тернарный»: \texttt{local v = cond and A or B}.
\item Диапазоны: индексы линейки \texttt{0..28}; буфер \texttt{Ledbar.new(29)}.
\end{itemize}
